\section{Materials and methods}
\label{sec:methods}

Figure~\ref{fig:workflow} summarizes the workflow: compilation and
cleaning of the literature data, the two-task evaluation protocol,
leakage-controlled model pipelines, and the uncertainty and
interpretation layers.

\begin{figure}[t]
\centering
\includegraphics[width=\linewidth]{fig_workflow}
\caption{Evaluation workflow}
\label{fig:workflow}
\end{figure}

\subsection{Data curation}
\label{sec:data}

We compiled 1,342 batch-adsorption experiments for LDH--micropollutant
systems from 59 peer-reviewed studies. Each record holds 30 descriptors in
four groups: micropollutant properties (compound identity, log
$K_{\mathrm{ow}}$, molecular weight, the Abraham solvation parameters
$E$, $S$, $A$, $B$, $V$, and topological polar surface area), LDH
composition and synthesis (divalent and trivalent metal, molar ratio,
synthesis route, hydrothermal temperature and time, calcination
conditions, and composition-averaged elemental descriptors), material
texture (interlayer spacing $d_{003}$, BET surface area, pore volume,
average pore diameter, and point of zero charge pH$_{\mathrm{pzc}}$), and
operating conditions (pH, temperature, contact time, adsorbent dose, and
initial concentration). Table~\ref{tab:stats} summarizes the numeric
descriptors; the prediction target is the removal ratio (RR\%).

\textbf{Cleaning.} Values stored as fractions were converted to percent by
a range-aware rule that leaves already-percent data untouched, and
category labels were normalized (e.g., merging spelling variants of
\emph{co-precipitation}). Records identical in every descriptor were
reduced to a single occurrence, since duplicates spanning a train--test
boundary would let the test be answered by lookup; deduplication removed
83 records. Missing texture values were imputed by a
$k$-nearest-neighbor imputer computed in standardized space and mapped
back to physical units; the imputer, like every fitted transform in this
work, was trained on training folds only. The source study of each
record is retained as a grouping key and is never used as a model input.
Details of the cleaning rules are in Text~S1.

\begin{table}[t]
\centering
\caption{Summary statistics of the numeric descriptors}
\label{tab:stats}
\begin{tabular}{lccc}
\toprule
Descriptor & Median & Range & Missing (\%) \\
\midrule
\multicolumn{4}{l}{\emph{Micropollutant}} \\
log $K_{\mathrm{ow}}$ & $-0.02$ & $-2.67$--$4.51$ & 0 \\
MW (g/mol) & 331 & 152--734 & 0 \\
$E$ / $S$ / $A$ / $B$ / $V$ & 1.81 / 2.10 / 0.70 / 1.70 / 2.30 & -- & 0 \\
TPSA (\AA$^2$) & 76.2 & 37.3--194.0 & 0 \\
\midrule
\multicolumn{4}{l}{\emph{Composition and synthesis}} \\
M(II)/M(III) ratio & 3.0 & 0.5--8.0 & 0 \\
Avg.\ atomic weight (g/mol) & 48.1 & 24.8--73.7 & 0 \\
Avg.\ d-electrons & 4.25 & 0--8.89 & 0 \\
Avg.\ Pauling electronegativity & 1.59 & 1.30--1.89 & 0 \\
Hydrothermal temp.\ (\textcelsius) & 70 & 25--160 & 13.3 \\
Hydrothermal time (h) & 24 & 12--40 & 13.3 \\
Calcination temp.\ (\textcelsius) & 400 & 200--550 & 65.9 \\
Calcination time (h) & 4 & 1--12 & 65.9 \\
\midrule
\multicolumn{4}{l}{\emph{Material texture}} \\
$d_{003}$ (nm) & 0.77 & 0.25--1.13 & 2.5 \\
BET (m$^2$/g) & 82.0 & 0.4--225.7 & 7.9 \\
PV (cm$^3$/g) & 0.10 & 0--1.46 & 25.3 \\
APD (nm) & 8.4 & 0.2--64.2 & 13.9 \\
pH$_{\mathrm{pzc}}$ & 7.5 & 5.5--11.6 & 24.6 \\
\midrule
\multicolumn{4}{l}{\emph{Operating conditions}} \\
pH & 6.7 & 2.0--12.0 & 0 \\
Temperature (\textcelsius) & 25 & 10--55 & 0 \\
Contact time (h) & 1.5 & 0.02--192 & 0 \\
Dose (g/L) & 0.9 & 0.01--50 & 0 \\
$C_i$ (mg/L) & 40 & 1--1688 & 0 \\
\bottomrule
\end{tabular}
\end{table}

\subsection{Evaluation protocol}
\label{sec:protocol}

\textbf{Two prediction tasks.} We evaluate every model under two
protocols that correspond to the two services expected from
literature-trained models. Task (T1), \emph{within-system interpolation},
asks how well a model predicts RR\% for a studied
adsorbent--pollutant system at new operating conditions. Task (T2),
\emph{cross-study generalization}, asks how well the model predicts
systems it has never seen. Intuitively, (T1) allows the model to have
seen the material and the pollutant at other conditions, while (T2) does
not.

\textbf{Repeated splits for (T1).} A single random split can be lucky or
unlucky, so (T1) is measured over 20 seeded 80/20 record-level splits of
the duplicate-free data, and scores are reported as mean $\pm$ standard
deviation across repeats. Hyperparameters are tuned once (Section
\ref{sec:models}) and held fixed across repeats, so the spread reflects
sampling variability rather than tuning noise.

\textbf{Leave-one-study-out audit for (T2).} Each of the 59 studies $s$
is held out once; the model $f_{-s}$ is refit on the remaining 58
studies with fixed hyperparameters and scored on the held-out records:
\begin{equation}
\mathrm{MAE}_s \;=\; \frac{1}{n_s}\sum_{i \in s}
\bigl|\, y_i - f_{-s}(\mathbf{x}_i) \,\bigr| .
\label{eq:loso}
\end{equation}
Reporting the distribution of $\mathrm{MAE}_s$ across studies, rather
than one pooled score, exposes which systems transfer and which do not.
A learning curve over the number of training studies tests whether
cross-study error is limited by data volume or by data diversity.

\textbf{Naive baselines.} To make the error magnitudes interpretable, we
score three predictors that use no descriptors: the global training-mean;
the per-pollutant mean (the average RR\% of the same compound in the
training data, falling back to the global mean for unseen compounds); and,
for (T1) only, the per-study mean. A model is useful for screening only
if it beats these lookups on unseen studies.

\textbf{Leakage control.} Imputation, scaling, and categorical encoding
are embedded in each model pipeline and fitted on training folds only.
Hyperparameter tuning uses cross-validation folds that differ from the
folds used for reporting, and grouped tasks use group-aware folds
throughout. These controls address the leakage mechanisms documented for
applied ML \citep{kapoor2023leakage}, which we found to inflate reported
skill substantially when absent (Section~\ref{sec:rq2}).

\textbf{Deployment workflow.} The protocol doubles as a reusable recipe
for future compilations (Fig.~\ref{fig:workflow}). Step~1: curate and
normalize the literature records, retaining the source-study identifier.
Step~2: remove exact duplicates and declare the prediction task:
interpolation within studied systems or screening of new ones.
Step~3: split accordingly (record-level for the former, study-grouped
for the latter) and fit pipelines whose preprocessing sees training
folds only. Step~4: wrap the chosen model in split-conformal
calibration matched to the declared task and report coverage
empirically. Every step is implemented in the released pipeline
(\texttt{ldh\_pipeline.py}), so re-running the audit on an extended
compilation requires only the new spreadsheet.

\subsection{Regression models}
\label{sec:models}

We compare seven regressors spanning the model families used in the
adsorption-ML literature, chosen so that their inductive biases differ.
Throughout, $\mathbf{x}_i \in \mathbb{R}^{p}$ denotes the preprocessed
descriptor vector of record $i$ and $\hat{y}_i$ its predicted RR\%.

\textbf{Linear regression (LR)} fits one additive coefficient per
descriptor,
\begin{equation}
\hat{y} = \beta_0 + \sum_{j=1}^{p} \beta_j x_j ,
\label{eq:lr}
\end{equation}
with $\boldsymbol{\beta}$ estimated by least squares. It serves as the
interpretable floor: any nonlinearity or interaction in the adsorption
response appears as a gap between LR and the flexible models.

\textbf{$k$-nearest neighbors (KNN)} predicts the average target of the
$k$ closest training records in standardized feature space
\citep{cover1967knn},
\begin{equation}
\hat{y}(\mathbf{x}) = \frac{1}{k}\sum_{i \in \mathcal{N}_k(\mathbf{x})} y_i ,
\label{eq:knn}
\end{equation}
where $\mathcal{N}_k(\mathbf{x})$ holds the $k$ nearest neighbors under
Euclidean distance and $k$ is grid-searched between 1 and 9. KNN has no
functional form at all; its performance measures how far pure similarity
lookup carries prediction, which makes it a diagnostic for
near-duplicate structure in the data.

\textbf{Artificial neural network (ANN)} is a single-hidden-layer
perceptron trained by Adam \citep{rumelhart1986backprop},
\begin{equation}
\hat{y} = \mathbf{w}^{\top} \sigma\!\left(W\mathbf{x} + \mathbf{b}\right) + b_0 ,
\label{eq:ann}
\end{equation}
with 32 hidden units and rectified-linear activation
$\sigma(z) = \max(0, z)$, matching the compact architectures used in
prior adsorption studies. Larger architectures were not explored because
the dataset holds fewer than $10^3$ effective degrees of freedom per
feature.

\textbf{Cubist} builds a set of rules $R_m$ whose leaves hold linear
models \citep{quinlan1992m5},
\begin{equation}
\hat{y}(\mathbf{x}) = \operatorname*{avg}_{m:\,\mathbf{x} \in R_m}
\Bigl( \beta_{0}^{(m)} + \textstyle\sum_j \beta_{j}^{(m)} x_j \Bigr),
\label{eq:cubist}
\end{equation}
optionally smoothed toward the targets of nearby training instances;
committees and neighbors are grid-searched. Unlike constant-leaf trees,
its leaf-linear structure can extrapolate moderately beyond the training
range of a rule, a property that becomes relevant under (T2).

\textbf{Random forest (RF)} averages $B$ decorrelated trees grown on
bootstrap samples \citep{breiman2001rf},
\begin{equation}
\hat{y}(\mathbf{x}) = \frac{1}{B}\sum_{b=1}^{B} T_b(\mathbf{x}),
\label{eq:rf}
\end{equation}
providing a variance-reduced piecewise-constant baseline among the
ensembles.

\textbf{Gradient-boosted trees} build the prediction stagewise, each new
tree fitting the residual structure left by its predecessors,
\begin{equation}
F_m(\mathbf{x}) = F_{m-1}(\mathbf{x}) + \nu\, h_m(\mathbf{x}),
\label{eq:gb}
\end{equation}
with learning rate $\nu$ and base trees $h_m$. \textbf{XGBoost}
implements Eq.~(\ref{eq:gb}) with second-order updates and a regularized
objective that penalizes the number of leaves and their weights
\citep{chen2016xgboost}; \textbf{histogram gradient boosting (HistGB)}
is the scikit-learn implementation with histogram-binned split finding
\citep{ke2017lightgbm, pedregosa2011sklearn}. Both are the strongest
current defaults for tabular data.

Tree-ensemble hyperparameters were tuned with the TPE sampler of Optuna
over 30 trials \citep{akiba2019optuna} on cross-validation folds of the
training data. Scale-sensitive models (LR, KNN, ANN, Cubist) receive a
standardizer inside their pipeline; tree ensembles operate on raw feature
scales. Search spaces and selected values are in Table~S1; all splits,
folds, and samplers are seeded for reproducibility.

\subsection{Feature screening and selection}
\label{sec:screening}

Multicollinearity is diagnosed with the variance inflation factor,
\begin{equation}
\mathrm{VIF}_j = \frac{1}{1 - R_j^{2}},
\label{eq:vif}
\end{equation}
where $R_j^{2}$ is the coefficient of determination of descriptor $j$
regressed on all other descriptors (targets excluded). Because
descriptors are constant within a study, the study-cluster structure
drives VIF values orders of magnitude beyond the usual threshold of 10
(Section~\ref{sec:eda}), so VIF is used here as a redundancy diagnostic
rather than an automatic elimination rule.

Feature selection is then evaluated the same way as the models
themselves: three nested descriptor sets (the full set of 30, a
correlation-screened set of 21 that drops the collinear Abraham and
composition-averaged block, and the 10 descriptors with the best median
permutation-importance rank) are compared under both tasks with fixed
hyperparameters. Judging a feature set only under (T1) would inherit
the same optimism as judging a model there, so both tasks are reported.

\subsection{Uncertainty and applicability domain}
\label{sec:uq}

\textbf{Split-conformal intervals.} We wrap the best model in
split-conformal prediction \citep{vovk2005conformal, lei2018conformal}.
A calibration subset of $n$ records is withheld from training, absolute
residuals $r_i = |y_i - \hat{y}_i|$ are computed on it, and the interval
half-width is their conformal quantile
\begin{equation}
q \;=\; \mathrm{Quantile}\!\left(\{r_i\}_{i=1}^{n};\;
\frac{\lceil (n+1)(1-\alpha) \rceil}{n}\right),
\label{eq:conformal}
\end{equation}
so that each test prediction is reported as $[\hat{y}-q,\ \hat{y}+q]$
with nominal coverage $1-\alpha = 0.9$. For task (T2) the calibration
subset is drawn as whole studies, so the calibration residuals reflect
the new-study setting. Conformal validity assumes exchangeability
between calibration and test units, which grouped data satisfy only
approximately; we therefore report empirical coverage for both tasks
rather than assuming it.

\textbf{Applicability domain.} As a candidate trust flag we compute, for
each test record, the mean standardized distance to its five nearest
training records in the preprocessed feature space. Records far from the
training distribution should receive predictions with caution; we
validate the flag by correlating it with realized absolute error rather
than assuming it works.

\subsection{Feature engineering and attribution}
\label{sec:features}

Beyond the raw descriptors, we construct three features with mechanistic
meaning: the electrostatic driving force pH$-$pH$_{\mathrm{pzc}}$, whose
sign determines whether the LDH surface is positively or negatively
charged relative to the solution; the initial loading ratio $C_i$/dose;
and the interaction log $K_{\mathrm{ow}} \times$
(pH$-$pH$_{\mathrm{pzc}}$) coupling hydrophobicity with surface charge.
Both tasks are re-evaluated with and without these features. Feature
attributions are computed with SHAP \citep{lundberg2017shap} on the
held-out test set using the fitted models, complemented by permutation
importance ranks across all seven models and one-dimensional partial
dependence profiles, and the resulting patterns are read against the
expected electrostatic and hydrophobic mechanisms.

\subsection{Performance metrics}
\label{sec:metrics}

Model skill is reported as the coefficient of determination,
\begin{equation}
R^2 = 1 - \frac{\sum_i (y_i - \hat{y}_i)^2}{\sum_i (y_i - \bar{y})^2},
\label{eq:r2}
\end{equation}
the root mean square error,
$\mathrm{RMSE} = \smash{\sqrt{\tfrac{1}{n}\sum_i (y_i - \hat{y}_i)^2}}$,
and the mean absolute error,
$\mathrm{MAE} = \tfrac{1}{n}\sum_i |y_i - \hat{y}_i|$,
each computed on three bases with distinct meanings: the training
records (in-sample fit), cross-validation folds of the training set
(validation), and the held-out test records of the task at hand. All
three are reported to keep the in-sample ceiling and the generalization
estimate visibly separate. Wall-clock training and prediction times are
reported in Section~\ref{sec:timing}.
